\documentclass[12pt,a4paper]{article}

% ── Packages ──
\usepackage[margin=2.5cm]{geometry}
\usepackage{fontenc}
\usepackage[utf8]{inputenc}
\usepackage{xcolor}
\usepackage{enumitem}
\usepackage{titlesec}
\usepackage{setspace}
\usepackage{microtype}
\usepackage{hyperref}
\usepackage{fancyhdr}
\usepackage{graphicx}

% ── Colors ──
\definecolor{huaweired}{HTML}{CF0A2C}
\definecolor{aiviolet}{HTML}{8E44AD}
\definecolor{ossblue}{HTML}{2E86AB}
\definecolor{bssgreen}{HTML}{2ECC71}
\definecolor{cvmorange}{HTML}{E67E22}
\definecolor{darkbg}{HTML}{1A1A2E}
\definecolor{stagedir}{HTML}{888888}

% ── Formatting ──
\titleformat{\section}{\Large\bfseries\color{huaweired}}{}{0em}{}[\vspace{-0.5em}\color{huaweired}\rule{\textwidth}{1.5pt}]
\titleformat{\subsection}{\large\bfseries\color{aiviolet}}{}{0em}{}

\setstretch{1.4}
\setlength{\parindent}{0pt}
\setlength{\parskip}{0.6em}

\pagestyle{fancy}
\fancyhf{}
\fancyhead[L]{\small\color{stagedir}\textit{Speech Script --- Telecom
        NeXoligence}}
\fancyhead[R]{\small\color{stagedir}\textit{Souhayl Guenichi}}
\fancyfoot[C]{\small\color{stagedir}\thepage}
\renewcommand{\headrulewidth}{0.5pt}

% ── Stage Direction Command ──
\newcommand{\stage}[1]{\par\textit{\color{stagedir}[#1]}\par}
\newcommand{\beat}{\par\vspace{0.4em}}
\newcommand{\pause}{\par\textit{\color{stagedir}[pause --- let it sink
            in]}\par}
\newcommand{\keyword}[1]{\textbf{\color{aiviolet}#1}}
\newcommand{\emphasis}[1]{\textbf{\color{huaweired}#1}}

\begin{document}

% ═══════════════════════════════════════════════════════
% TITLE PAGE
% ═══════════════════════════════════════════════════════
\begin{titlepage}
    \centering
    \vspace*{3cm}

    {\Huge\bfseries\color{huaweired} SPEECH SCRIPT}\\[0.5cm]
    {\LARGE\color{aiviolet} Supervisor Meeting --- Progress Overview}\\[0.3cm]
    {\Large\color{ossblue} Cloud-Native AI Operations Agent for CEM--CVM
    Intelligence}\\[2cm]

    {\large\color{darkbg} Souhayl Guenichi}\\[0.3cm]
    {\normalsize ESPRIT $\times$ Huawei Tunisia}\\[0.3cm]
    {\normalsize Supervised by: Mrs.\ Mariem Bouzouita (ESPRIT) \& Mr.\ Jackie Chen
    (Huawei)}\\[2cm]

    \textit{\color{stagedir}Script for the supervisor progress meeting.\\
        This is an overview presentation --- not the final defense.\\
        Focus on the idea, approach, outputs, and next steps.\\
        No code, no ports, no implementation details.\\
        Target: 15--20 minutes total.}

    \vfill
    {\small March 2026}
\end{titlepage}

\tableofcontents
\newpage

% ═══════════════════════════════════════════════════════
% SLIDE 1 — TITLE (30 seconds)
% ═══════════════════════════════════════════════════════
\section{SLIDE 1 --- TITLE (30 seconds)}

\stage{Title slide on screen. Stand confidently.}

``Good morning everyone. Thank you for your time today.''

\beat

``My name is Souhayl Guenichi, fifth-year engineering student at ESPRIT. I am
doing my PFE at \emphasis{Huawei Tunisia}, in the Cloud IT / Sales-Solution
team.''

\beat

``Today I want to walk you through the project I have been building: a
\keyword{Cloud-Native AI Operations Agent for CEM--CVM Intelligence}. In simple
terms --- an autonomous AI system for telecom that \emphasis{detects problems,
    decides what to do, and acts on it}.''

\beat

``This is a progress meeting, so I will focus on the big picture --- the idea,
the approach, what we have built so far, and what is next.''

% ═══════════════════════════════════════════════════════
% SLIDE 2 — THE PROBLEM (2 minutes)
% ═══════════════════════════════════════════════════════
\newpage
\section{SLIDE 2 --- THE PROBLEM (2 minutes)}

\stage{Show Problem slide --- OSS silo, BSS silo, red GAP in the middle.}

``Let me start with the problem.''

\beat

``Telecom operators have two data worlds that \emphasis{never talk to each
    other}. On one side, the \keyword{network} --- OSS data: throughput, latency,
packet loss, signal quality. On the other side, the \keyword{business} --- BSS
data: revenue, churn, data usage, ARPU.''

\beat

``When the network degrades in a region, the business team sees the revenue
impact \emphasis{days later}. Nobody connects the dots. And more importantly
--- \emphasis{nobody takes action}. There is no system that bridges this gap
and acts autonomously.''

\pause

``That is what we set out to build.''

% ═══════════════════════════════════════════════════════
% SLIDE 3 — THE SOLUTION (3 minutes)
% ═══════════════════════════════════════════════════════
\newpage
\section{SLIDE 3 --- THE SOLUTION (3 minutes)}

\stage{Show CEM → Agent → CVM diagram. This is the key visual.}

``Our solution is an \keyword{AI Operations Agent} that sits right in the
middle --- between Huawei's \keyword{CEM} platform, SmartCare, and the
operator's \keyword{CVM} system.''

\beat

``The agent has three core capabilities:''

\beat

``First, \emphasis{DETECT} --- it uses machine learning to detect anomalies in
the network and predict SLA risk.''

\beat

``Second, \emphasis{DECIDE} --- it correlates network metrics with business
metrics. When the network degrades, what happens to revenue? It quantifies that
relationship.''

\beat

``Third, \emphasis{ACT} --- and this is the key differentiator --- it does not
just give you a report. It \keyword{generates autonomous action
    recommendations}: alerts, churn prevention campaigns, upsell triggers,
escalation tickets.''

\pause

``This positions us at \keyword{ADN Level 4} --- Highly Autonomous --- in
Huawei's Autonomous Driving Network framework. At L3, the system recommends and
humans decide. At L4, the agent \emphasis{decides and acts}. Humans only handle
edge cases.''

\beat

``And we work with \emphasis{real data from Tunisie Telecom} --- not synthetic,
not simulated. Real network KPIs, real subscriber records, anonymised.''

% ═══════════════════════════════════════════════════════
% SLIDE 4 — THE PIPELINE (2 minutes)
% ═══════════════════════════════════════════════════════
\newpage
\section{SLIDE 4 --- HOW IT WORKS: 5-PHASE PIPELINE (2 minutes)}

\stage{Show 5-Phase Pipeline diagram. Walk through each phase briefly.}

``Internally, the agent runs a \keyword{5-phase pipeline}:''

\beat

``\keyword{Phase 1} --- we ingest real operator data from Tunisie Telecom into
our data lake.''

\beat

``\keyword{Phase 2} --- we engineer meaningful features from the raw network
and business data.''

\beat

``\keyword{Phase 3} --- we run AI models: predict SLA risk, detect network
anomalies, detect revenue anomalies, and find correlations between the two
domains.''

\beat

``\keyword{Phase 4} --- and this is what makes us L4 --- the \keyword{Decision
    and Action Engine}. It evaluates the ML outputs against business rules and
\emphasis{generates typed recommendations with priority levels}.''

\beat

``\keyword{Phase 5} --- everything is persisted: curated data plus actions,
ready for CVM consumption.''

\beat

``Each phase is color-coded --- teal for ingestion, blue for engineering,
violet for AI, red for actions, green for persistence. The whole pipeline runs
end-to-end with one command.''

% ═══════════════════════════════════════════════════════
% SLIDE 5 — AI MODELS (2 minutes)
% ═══════════════════════════════════════════════════════
\newpage
\section{SLIDE 5 --- THE AI MODELS (2 minutes)}

\stage{Show ML Pipeline with Action Engine diagram.}

``The intelligence layer has four AI components:''

\beat

``\keyword{SLA Risk Prediction} --- a Gradient Boosting model that answers: how
likely is this region to breach its SLA? It gives a score from 0 to 1 based on
9 network features.''

\beat

``\keyword{Network Anomaly Detection} --- an Isolation Forest model that flags
cells behaving abnormally.''

\beat

``\keyword{Revenue Anomaly Detection} --- same algorithm, different domain. It
catches unexpected drops or spikes in subscriber revenue.''

\beat

``\keyword{O+B Correlation Engine} --- this is where the cross-domain magic
happens. We compute Pearson and Spearman correlations between network metrics
and business metrics. When latency goes up, does revenue go down? Now we can
\emphasis{quantify} it.''

\beat

``All these outputs feed into the \keyword{Action Engine}, which decides what
to do about them.''

% ═══════════════════════════════════════════════════════
% SLIDE 6 — THE ACTION ENGINE (2 minutes)
% ═══════════════════════════════════════════════════════
\newpage
\section{SLIDE 6 --- THE ACTION ENGINE (2 minutes)}

\stage{Show Action Engine Rules diagram. Speak with emphasis --- this is the L4
    differentiator.}

``This slide is the heart of the project. The \keyword{Decision and Action
    Engine} is what makes us an \emphasis{autonomous agent}, not just an analytics
dashboard.''

\beat

``The engine takes ML outputs and evaluates them against business rules:''

\beat

``If SLA risk is above 0.7 --- \emphasis{proactive alert}, critical priority.
The network is about to breach, act now.''

\beat

``If there is an OSS anomaly correlated with a revenue drop --- \emphasis{churn
    prevention campaign}. Subscribers in that region are at risk.''

\beat

``If the network is healthy but revenue spikes --- \emphasis{upsell
    opportunity}. These are power users, offer them more.''

\beat

``If there is a significant O+B correlation --- \emphasis{insight report} for
the CVM strategy team.''

\beat

``If multiple anomalies hit the same region --- \emphasis{escalation ticket},
critical. Something serious is happening.''

\pause

``Each action has a lifecycle: generated, pending, approved or dismissed. The
CVM system can consume these through our API. The agent decides and acts.
Humans handle edge cases only. \emphasis{That is L4.}''

% ═══════════════════════════════════════════════════════
% SLIDE 7 — ARCHITECTURE (1.5 minutes)
% ═══════════════════════════════════════════════════════
\newpage
\section{SLIDE 7 --- ARCHITECTURE (1.5 minutes)}

\stage{Show Container Architecture diagram + HCS Mapping.}

``A quick look at the architecture. The system runs as \keyword{5 containerised
    microservices} in Docker Compose:''

\beat

``An API Gateway for consuming results. An AI Service for ML inference. The
Pipeline Worker that orchestrates the 5 phases. PostgreSQL for our 8-table
structured store. And MinIO as a three-layer data lake --- raw, processed,
curated.''

\beat

``The entire stack runs locally with one command, and it is designed for
\keyword{cloud portability}. The same Docker images deploy directly to Huawei
Cloud Stack --- containers to ECS, MinIO to OBS, PostgreSQL to RDS. Same code,
different infrastructure.''

% ═══════════════════════════════════════════════════════
% SLIDE 8 — DATA (1.5 minutes)
% ═══════════════════════════════════════════════════════
\newpage
\section{SLIDE 8 --- DATA (1.5 minutes)}

\stage{Show Data slide.}

``A word on data, because this is a major differentiator.''

\beat

``Most PFE projects use synthetic data. We work with \emphasis{real anonymised
    production data from Tunisie Telecom}. Real network KPIs across regions. Real
subscriber profiles with revenue and usage patterns.''

\beat

``Everything is anonymised --- hashed identifiers, no personal information,
\mbox{gouvernorat-level} location only.''

\beat

``We also have a dual-mode design: real data is the default, with a synthetic
fallback for demos and testing. The system works from day one and only gets
better with real data.''

\beat

``Our database has 8 tables with full traceability --- every anomaly, every
risk score, every action is linked back to the pipeline run that produced it.''

% ═══════════════════════════════════════════════════════
% SLIDE 9 — COMPETITIVE POSITIONING (1.5 minutes)
% ═══════════════════════════════════════════════════════
\newpage
\section{SLIDE 9 --- WHAT MAKES US DIFFERENT (1.5 minutes)}

\stage{Show Competitive Comparison. Speak with conviction.}

``How do we compare to what exists today?''

\beat

``SELFNET does self-healing 5G --- no real data, no O+B convergence, no
actions.''

\beat

``ETSI ZSM defines zero-touch management --- but it is a specification, not an
implementation.''

\beat

``Nokia AVA does telecom analytics with real data --- but no O+B bridge and no
autonomous actions.''

\beat

``Huawei SmartCare does CEM with partial O+B --- but still no autonomous
actions.''

\pause

``Our project is the \emphasis{only one} combining all four: \keyword{real
    data, O+B convergence, cloud-native architecture, and autonomous actions}. That
is our contribution.''

% ═══════════════════════════════════════════════════════
% SLIDE 10 — DELIVERABLES (1 minute)
% ═══════════════════════════════════════════════════════
\newpage
\section{SLIDE 10 --- OUTPUTS SO FAR (1 minute)}

\stage{Show Deliverables slide. Go through quickly.}

``Here is what we have delivered so far:''

\beat

``A \keyword{working Docker stack} --- 5 services, runs with one command. Three
\keyword{trained ML models} on real Tunisie Telecom data. A \keyword{Decision
    and Action Engine} generating typed recommendations. A complete
\keyword{three-layer data lake}. A \keyword{REST API} for consuming results and
managing actions. An \keyword{8-table PostgreSQL schema} with full pipeline
lineage. The \keyword{Huawei Cloud Stack architecture} design. And
\keyword{technical documentation} including the report and architecture docs.''

\beat

``All of this is functional and running today.''

% ═══════════════════════════════════════════════════════
% SLIDE 11 — NEXT STEPS & CLOSE (1 minute)
% ═══════════════════════════════════════════════════════
\newpage
\section{SLIDE 11 --- NEXT STEPS \& CLOSE (1 minute)}

\stage{Show Next Steps slide. Slow down for the close.}

``Looking ahead:''

\beat

``\emphasis{Immediately} --- finalise the action engine rules with feedback
from the Huawei team.''

\beat

``\emphasis{Short-term} --- deploy to the actual Huawei Cloud Stack
environment.''

\beat

``\emphasis{Medium-term} --- add real-time streaming and more advanced ML
models like LSTM and transformers.''

\beat

``\emphasis{Long-term} --- reinforcement learning for action optimisation,
moving toward ADN Level 5.''

\pause

\stage{Make eye contact. Smile.}

``That is where we are. I would love to hear your feedback and questions.''

\beat

``Thank you.''

% ═══════════════════════════════════════════════════════
% APPENDIX — QUICK CHEAT SHEET
% ═══════════════════════════════════════════════════════
\newpage
\section{APPENDIX --- Quick Cheat Sheet for Q\&A}

\subsection{Likely Questions \& Answers}

\textbf{Q: Why these specific ML models?}\\
A: GBR gives interpretable feature importance for SLA risk. Isolation Forest is
ideal for unsupervised anomaly detection without labeled data. Pearson/Spearman
quantifies cross-domain correlations. Simple, explainable, effective.

\textbf{Q: What makes this L4 and not L3?}\\
A: At L3 the system recommends and humans approve everything. At L4, the agent
generates actions autonomously --- alerts, campaigns, escalations. Humans only
intervene for critical edge cases. The Action Engine is what makes the
difference.

\textbf{Q: Why not go straight to L5?}\\
A: L5 means fully autonomous, zero human involvement. Unrealistic for a PFE
scope. L4 is honest and defensible --- autonomous for routine cases, human
oversight for exceptions.

\textbf{Q: Is this just a dashboard?}\\
A: No. A dashboard displays data. Our agent \emphasis{takes action}. It
generates typed recommendations with confidence and priority. CVM systems
consume these via API and act on them. That is the fundamental difference.

\textbf{Q: How is the data anonymised?}\\
A: Subscriber IDs are SHA-256 hashed with salt. Cell identifiers pseudonymised.
Gouvernorat-level location only. No PII in the processed or curated layers.

\textbf{Q: Why cloud-native?}\\
A: Huawei's strategy is cloud. Docker Compose gives us portability --- same
images deploy to HCS (ECS, OBS, RDS). We prove it works locally and the path to
production is clear.

\textbf{Q: What if you do not get more real data?}\\
A: Dual-mode design. Synthetic fallback is built in. The system trains and runs
on synthetic, and improves when real data arrives. We already have initial real
data from TT.

\textbf{Q: What is the O+B convergence concretely?}\\
A: Quantifying the statistical relationship between network metrics (latency,
throughput) and business metrics (revenue, churn). Example: when latency
increases by 20ms, revenue drops by X\%. No existing tool does this
automatically.

\textbf{Q: What is the next milestone?}\\
A: Deploy to Huawei Cloud Stack and validate the action engine rules with the
Huawei operations team.

\end{document}
